problem:  
Let \( (G/H,g) \) be a connected, simply connected geodesic orbit (g.o.) Riemannian homogeneous space with reductive decomposition \( \mathfrak{g} = \mathfrak{h} \oplus \mathfrak{m} \), and \(g\) induced by an \( \mathrm{Ad}(H)\)-invariant inner product on \(\mathfrak{m}\).  
For each **geodesic vector** \(X\in \mathfrak{g}\), denote its tangent component by \(X_\mathfrak{m}\in \mathfrak{m}\), and define the Jacobi operator  
\[
J_{X_\mathfrak{m}} := R(X_\mathfrak{m},\cdot)X_\mathfrak{m} : \mathfrak{m}\to\mathfrak{m},
\]
where \(R\) is the curvature tensor at the base point \(o=eH\).  
Assume that, for every such geodesic vector \(X\), the spectrum of \(J_{X_\mathfrak{m}}\) (counted with multiplicities) is constant along the homogeneous geodesic \(\gamma_X(t)=\exp(tX)\cdot o\); i.e. the eigenvalues of \(J_{X_\mathfrak{m}}\) do not depend on \(t\).  
Conjecture: Does this condition imply that \( (G/H,g) \) is naturally reductive? If not, can one construct a non–naturally–reductive example satisfying this constant‑Jacobi‑spectrum property?  

Why is it a "good" problem:  
This version is well‑defined and formally correct, projecting geodesic vectors properly and formulating the spectral condition precisely. It explores a new rigidity property arising from the **spectral constancy of the Jacobi operator along homogeneous geodesics**, a condition lying halfway between the full curvature parallelism of symmetric spaces (\(\nabla R=0\)) and the algebraic symmetry of natural reductivity.  
The problem directly connects the **Lie‑algebraic structure** of a homogeneous space with the **analytic behavior of curvature along geodesics**, bridging homogeneous Riemannian geometry, representation theory, and spectral geometry. The question—whether spectral invariance forces natural reductivity—is both simple in statement and rich in implication: a positive answer would reveal a spectral rigidity principle for naturally reductive spaces, while a negative one would yield new, unexplored homogeneous geometries defined by a subtle curvature condition not yet classified.  
Its novelty lies in importing ideas from **Osserman‑type spectral geometry** into the context of **homogeneous geodesic orbit spaces**, creating a cross‑fertilized research direction that remains uncharted. This combination of clean formulation, broad interfield relevance, and potential to unveil new rigidity phenomena makes it a genuinely "good" research‑level problem.